\section*{Chapter \ref{ch:iv1:estimation} --- Estimation}
\markboth{\protect\omTocMark{}{Answers: Estimation}}{\protect\omTocMark{}{Answers: Estimation}}

\begin{solution}{iq:iv1:estimation:1}
A year has about $3.16 \times 10^7$ seconds, which is close to $\pi \times 10^7$. Trading hours hold about $5.9$ million of them.

The year is $365.25 \cdot 86\,400 = 31\,557\,600$ seconds. The shortcut $\pi \times 10^7$ is within half a percent, and it is worth memorising. A trading year is $252 \cdot 6.5 \cdot 3600 = 5\,896\,800$ seconds, so the market is open for about $18.7\%$ of the time, a little less than a fifth.

\iqfollow{How many trading minutes in a year?}{$252 \cdot 390 = 98\,280$, or about $100\,000$.}
\iqwhy{Annualising intraday figures needs the trading clock. A rate of one event per trading second is about six million events a year, not thirty million.}
\end{solution}

\begin{solution}{iq:iv1:estimation:2}
About three months, with an interval from six weeks to six months.

The pool holds $50 \cdot 25 \cdot 2 = 2500$ cubic metres. A garden hose fills a 10-litre bucket in about 30 seconds, which is 20 litres a minute, or $1.2$ cubic metres an hour. So the pool takes $2500/1.2 \approx 2083$ hours, or 87 days running without pause.

The volume is exact, so the whole uncertainty lies in the flow. Between 10 and 40 litres a minute, a factor of two each way, the time runs from 174 days down to 43.

\iqwrong{``It would take a few days.'' Nobody has a feel for 2\,500 cubic metres. That is why you turn it into litres, and the hose into a bucket you can picture filling.}
\end{solution}

\begin{solution}{iq:iv1:estimation:3}
A 32-bit counter lasts about 72 minutes, a 64-bit one lasts about 585\,000 years, and a daily reset needs 35 bits.

The first is $2^{32}/10^6 \approx 4295$ seconds, using $2^{10} \approx 10^3$ to get $4 \times 10^9$ at once. The second is $2^{64}/10^6 \approx 1.8 \times 10^{13}$ seconds, and a year is about $\pi \times 10^7$ seconds. A day holds $6.5 \cdot 3600 \cdot 10^6 = 2.34 \times 10^{10}$ orders, which lies between $2^{34} \approx 1.7 \times 10^{10}$ and $2^{35} \approx 3.4 \times 10^{10}$. So 35 bits suffice.

\iqwhy{Thirty-two bits looks plentiful until a busy day. A counter that overflows at an hour and a quarter fails in production, not in a test with a small sample. Choose the width by this sum, and leave room for the rate to grow.}
\end{solution}

\begin{solution}{iq:iv1:estimation:4}
Not if days are independent: 5\% a day scales to about 79\% a year, and 20\% a year is only about $1.3\%$ a day.

With independent days, variances add, so volatility scales with the square root of time\index{square-root-of-time rule}. The shortcut is $\sqrt{252} \approx 16$. Either $5\% \cdot 16 \approx 80\%$ a year, or $20\%/16 \approx 1.25\%$ a day.

For both statements to hold, the daily moves would have to cancel almost completely. The actual yearly variance would be only $(20/79.4)^2 \approx 6.3\%$ of what independent days give, and that would need extreme mean reversion.

\iqfollow{Could a few violent days reconcile them?}{They can make the \emph{average} misleading. With 248 days at 1\% and 4 earnings days at 5\%, annual volatility is $\sqrt{248 + 4 \cdot 25} \approx 18.7\%$, so 5\% moves can occur. But then 5\% is not a typical day.}
\end{solution}

\begin{solution}{iq:iv1:estimation:5}
The fibre takes 6 milliseconds one way, a straight radio link takes 4, and a signal in fibre covers 20 centimetres per nanosecond.

Fibre runs at $200\,000$ km per second, so $1200/200\,000$ seconds is 6 ms. Radio through air takes $1200/300\,000$ seconds, or 4 ms, which saves up to 2 ms each way. In air the signal covers about 30 cm in a nanosecond.

A real fibre route is longer than the straight line, so the true saving is larger. A radio link is less reliable and carries less data, which is why both are used.

\iqwhy{This is the floor under any two-city trade: no software or hardware beats it. Within a building, 20 cm per nanosecond means a few metres of extra cable costs 10 to 20 nanoseconds, which can use up a hard-won software saving.}
\end{solution}

\begin{solution}{iq:iv1:estimation:6}
Power costs about \$657 a year, which is about a fifth of the server's yearly cost.

A year holds about 8\,760 hours, so $0.5$ kW draws $4380$ kWh, which costs \$438 at 10 cents. The building's overhead makes that $1.5 \cdot 438 = \$657$. The machine costs $10\,000/4 = \$2500$ a year, so power is $657/(2500 + 657) \approx 21\%$ of the total.

\iqwrong{Forgetting the overhead, or dividing by the machine cost alone. A share is a part over the whole, so the power belongs in both places.}
\iqwhy{When someone offers to save 10\% of a server's power, this tells you the saving is about 2\% of the bill, before any change in the rate of failure or rack space.}
\end{solution}

\begin{solution}{iq:iv1:estimation:7}
The cube has a side of about 22 metres, and it is worth about \$19 trillion.

Density is $19\,300$ kg per cubic metre, so $2 \times 10^8$ kg fills about $10\,360$ cubic metres. Since $22^3 = 10\,648$, the side is about $21.8$ metres. In ounces, $2 \times 10^{11}/31.1 \approx 6.43 \times 10^9$, which at \$3\,000 gives about \$$19.3$ trillion.

\iqfollow{Per person, with 8 billion people?}{About 25 grams each, which is less than one ounce.}
\iqwhy{The size is the surprise: all that value would fit in a modest building. That is a good habit in sizing, since a figure that sounds huge may be small once it is put into units that you can see.}
\end{solution}

\begin{solution}{iq:iv1:estimation:8}
About 2\,400, with an interval from 1\,000 to about 7\,000.

Work from demand\index{demand-side estimate}. One million people, each with about 6 cuts a year, means 6 million cuts. A hairdresser does about one cut every 45 minutes, so around 10 cuts in a day, on about 250 days: 2\,500 cuts a year. Dividing gives 2\,400.

For the interval, push both inputs together. At 4 cuts a year and 4\,000 cuts per worker you get 1\,000. At 10 cuts a year and 1\,500 cuts per worker you get about 6\,700. Their geometric mean\index{geometric mean of bounds} is about 2\,600, close to the central figure, which is a good sign that the two ends are balanced.

The check uses experience of a street, not of a haircut. The central figure means one worker for each 400 people. A neighbourhood of 2\,000 people would then have five workers, one or two salons, and that looks about right.

\iqwrong{Giving one number. The interviewer wants to see which inputs you are unsure of, and that you know the width of the answer.}
\end{solution}

\begin{solution}{iq:iv1:estimation:9}
About 650\,000, and the two routes give about 690\,000 and 600\,000. An interval of 500\,000 to 800\,000 is fair.

Seats and fittings take about a fifth of the $12 \cdot 2.5 \cdot 2 = 60$ cubic metres, which leaves 48.

Route one uses volume. A ball of radius $2.15$ cm has a volume $\tfrac43 \pi r^3 \approx 41.6$ cubic centimetres. Spheres poured at random fill about 60\% of a space, so the count is $48 \cdot 0.6/(41.6 \times 10^{-6}) \approx 692\,000$.

Route two stacks them in a cubic grid: every ball owns a cube of side $4.3$ cm, or $79.5$ cubic centimetres, and $48/(79.5 \times 10^{-6}) \approx 604\,000$. A grid fills only $\pi/6 \approx 52\%$ of the space, which is why this is lower.

\iqwrong{Dividing the volume of the bus by the volume of a ball and ignoring the gaps. That overstates it by more than half. Even the densest possible packing fills only 74\%, which caps the answer below a million.}
\end{solution}

\begin{solution}{iq:iv1:estimation:10}
About a million and a quarter, with an interval from about half a million to $2.5$ million. The two routes differ by about 20\%.

Route one uses passengers. With 8 billion people taking half a flight each a year on average, there are 4 billion passenger flights. At about $2.5$ hours each, that is $10^{10}$ passenger-hours in the $8760$ hours of a year, so $1.14$ million people are in the air at an average moment.

Route two uses aircraft. With about 25\,000 airliners, each flying 10 hours in 24, about $10\,400$ are aloft at any moment. At about 130 passengers each, that is $1.35$ million.

The agreement is what gives you confidence, because the routes share no input. The result means that about one person in 7\,000 is flying now. Averages hide a daily pattern, so the moment you ask may be somewhat above or below.

\iqwhy{Two independent routes reveal a mistaken input. If they had differed by a factor of ten, you would know that one of them is wrong, and that is more useful than a single confident number.}
\end{solution}

\begin{solution}{iq:iv1:estimation:11}
About 100\,000, with an interval from 50\,000 to 200\,000.

Start from demand. Of 300 million people, take 200 million cars. If each fills up about every 10 days, that is 20 million fills a day. A station with 8 pumps sees about 200 fills a day, or one every five minutes through a long day, so the country needs $20\,000\,000/200 = 100\,000$ stations. Between 400 and 100 fills a day, the count is 50\,000 to 200\,000.

Now check it from the supply side. That number means one station for every 3\,000 people. A station earns on each fill: at 40 litres and 5 cents of margin per litre, each fill brings \$2, or \$400 a day, or about \$146\,000 a year. That pays for a couple of workers on the site, and the shop adds more. So a station at 200 fills a day just about survives, and this suggests that the number cannot be much higher.

\iqwhy{Supply and demand give different kinds of check. The demand side tells you what is needed, and the supply side tells you what can stay in business, and a real business has to pass both.}
\end{solution}

\begin{solution}{iq:iv1:estimation:12}
From a quarter of your guess to four times it, when the factors are independent and the errors are normal in logarithms.

Take logarithms\index{lognormal errors}. A product of independent factors becomes a sum of independent errors, and in a sum variances add, so widths grow as the square root of the number of terms. A 90\% interval of a factor of 2 is $\pm \ln 2$ in logs, and so $\sigma = \ln 2/1.645 \approx 0.42$. With four factors the width doubles, to $\pm 2 \ln 2 = \pm \ln 4$. With nine it would triple, to a factor of 8.

Sixteen is what you get only if all four errors go the same way at once, as they do if the factors are perfectly correlated, because a single mistake drives them all. For independent factors, the interval from one sixteenth to sixteen would hold about 99.9\% of outcomes.

\iqfollow{Where is the mean of the product, compared with the central guess?}{Above it. A lognormal's mean is its median times $e^{\sigma^2/2}$, with $\sigma = 0.84$ here, so about $1.43$ times. A central Fermi figure is a median, and it understates an average.}
\iqwhy{Breaking a quantity into many pieces is what makes Fermi estimates work: independent errors partly cancel. But it works only if the pieces do not share one assumption.}
\end{solution}

\begin{solution}{iq:iv1:estimation:13}
About 1\,000 trades per second. At peak the feed is about 610\,000 messages per second, or about 200 megabits per second.

The daily traded value is $60 \times 10^{12} \cdot 2/252 \approx \$476$ billion. Dividing by \$20\,000 gives about $23.8$ million trades a day. A trading day has $6.5 \cdot 3600 = 23\,400$ seconds, so that is about 1\,020 trades a second.

With 30 messages a trade, the average is about $30\,500$ messages a second, and 20 times that is $610\,000$. At 40 bytes each that is $24.4$ million bytes a second, or $195$ megabits a second.

\iqwrong{Sizing a system for the average. Message rates cluster heavily around the open, the close and news, and a design that copes with the average falls behind exactly when the data matter most.}
\iqwhy{A figure of this kind tells you at once which hardware you need: a 1-gigabit link has room to spare here, but only until you add more venues.}
\end{solution}

\begin{solution}{iq:iv1:estimation:14}
About \$126 million a year, with an interval of \$63 million to \$252 million. The central figure means it keeps about a fifth of the half-spread.

Five percent of 10 billion is 500 million shares a day. At $0.1$ cents, which is the geometric middle of $0.05$ and $0.2$, each day brings \$500\,000, and over 252 days that is \$126 million. The ends scale directly with the cents kept.

The half-spread is $0.5$ cents, which is what a market maker earns if it trades at its quote and the price then stays put. Keeping $0.1$ cents means that about $0.4$ cents is lost to adverse selection\index{adverse selection}, fees and the cost of holding the position.

\iqwhy{Market-making revenue looks enormous in shares and thin per share. The spread is not the profit: most of it pays for being on the wrong side of informed traders.}
\end{solution}

\begin{solution}{iq:iv1:estimation:15}
About $11.8$ terabytes a year uncompressed, or about $2.4$ terabytes compressed. One uncompressed day takes about 47 seconds to read.

A day has $5000 \cdot 10 \cdot 6.5 \cdot 3600 = 1.17 \times 10^9$ updates. At 40 bytes each that is $46.8$ gigabytes. Over 252 days it comes to $11.8$ terabytes, and dividing by 5 gives $2.36$ terabytes. At one gigabyte a second, $46.8$ gigabytes takes about 47 seconds.

The average hides skew: a few hundred active names produce most of the updates, so a busy stock holds many times its $1/5000$ share of the day, about 9 milliseconds of reading, and a quiet one far less. Either way, reading one stock alone is far quicker than scanning the whole day, if the file layout lets you.

\iqwhy{The bottleneck in research is rarely the disk space for a year. It is the time to scan many days, and that is why tick stores partition by date and symbol.}
\end{solution}

\begin{solution}{iq:iv1:estimation:16}
About $2.7$ hours on one core and about 10 minutes on 16 cores, at which point the 6.5-minute disk read is close to binding.

The bars number $5000 \cdot 252 \cdot 20 \cdot 390 \approx 9.83 \times 10^9$. At 40 bytes each that is 393 gigabytes, so reading takes 393 seconds, or about $6.55$ minutes, whatever the number of cores. Processing at 1 microsecond a bar takes 9\,828 seconds, or $2.73$ hours, on one core.

The two run together, so on one core processing is by far the limit. On 16 cores the processing takes 614 seconds, still slower than the disk. They become equal at $9828/393 = 25$ cores, and past that more cores do nothing. You need a faster disk, or to read less data.

\iqwhy{The useful result is the ratio, not the $2.7$ hours. Know which resource binds before spending money: at 25 cores, more cores buy nothing, and storing the data by column, so that a run reads a tenth of it, helps everyone.}
\end{solution}

\begin{solution}{iq:iv1:estimation:17}
About \$$12.6$ billion a year, which is about $2.1$ basis points of the market's value.

Each trade has one side crossing the spread, and that side pays half of the quoted spread compared with the mid-price, 1 basis point here. So the daily toll is $500 \times 10^9 \cdot 10^{-4} = \$50$ million, and over 252 days that is \$$12.6$ billion. Against \$60 trillion, the total market value, it is $2.1$ basis points a year.

\iqwrong{Charging the full spread, which gives \$$25.2$ billion. A taker pays half the spread against the mid-price, and paying the full spread would need the same party to buy and then sell. Another mistake is to charge both sides of each trade: the resting side earns what the crossing side pays.}
\iqwhy{That is a useful ceiling on what spread-capturing businesses can share between them, before impact and adverse selection. It also shows how small the toll is for a long-term holder.}
\end{solution}

\begin{solution}{iq:iv1:estimation:18}
No. It implies 1\% of the stock's whole traded value every day, which is far more than any plausible edge.

The deciding figure is profit over volume: $\$1$ million over \$100 million is 100 basis points on every dollar traded by anyone that day, even if the strategy did all of the trading. A realistic share of 10\% means \$10 million traded, so it earns 10\% per dollar traded. That is five times the stock's daily standard deviation of 2\%, every day.

Even one day like that is a large outlier, and doing so every day with no losses cannot be skill in the ordinary sense. It suggests a stale mark, a data error, an unrealistic fill model or a misattribution.

\iqfollow{What would you ask next?}{How the fills were modelled, whether the P\&L was marked at prices at which you could really trade, and what the position size was at the close. These find the error faster than any argument about the signal.}
\iqwhy{Dividing profit by volume is the fastest check on a track record, and it needs no details of the strategy.}
\end{solution}

\begin{solution}{iq:iv1:estimation:19}
About $1.5$ molecules on average, so the chance of at least one is about 77\%.

First the atmosphere. Pressure is the weight of the air above a square metre, so each square metre carries $10^5/10 = 10^4$ kg. The Earth's surface is $4\pi R^2 \approx 5.15 \times 10^{14}$ square metres, so the atmosphere weighs about $5.15 \times 10^{18}$ kg. At 29 grams a mole that is about $1.77 \times 10^{20}$ moles, or about $1.06 \times 10^{44}$ molecules.

A breath is $0.5/24 \approx 0.0208$ moles, or $n \approx 1.25 \times 10^{22}$ molecules. Caesar's breath, mixed evenly, is a share $n/N$ of the atmosphere, and each of your $n$ molecules has that chance of being his, so the mean is $n^2/N \approx 1.46$.

The count is a sum of many very rare independent events, so it is Poisson\index{Poisson approximation}, and $P(\text{at least one}) = 1 - e^{-1.46} \approx 0.77$.

\iqwrong{Assuming the answer must be tiny, or vast, since one breath is so small and the air so large. The two huge numbers nearly cancel: $n$ appears squared, and $n^2$ happens to be close to $N$.}
\iqfollow{What assumption is weakest?}{Even mixing, and that the molecules survive unchanged. Over 2\,000 years air mixes well, but much of the carbon dioxide and some of the oxygen are taken up and replaced. Nitrogen is the most stable part, so the estimate is an upper end.}
\end{solution}

\begin{solution}{iq:iv1:estimation:20}
$52!$ is about $8 \times 10^{67}$, about $7 \times 10^{17}$ times the number of atoms in the Earth. The machine would have seen $4.4 \times 10^{26}$ orderings, and an even chance of a match needs about $10^{34}$ shuffles.

Use Stirling's formula\index{Stirling's formula}: $\ln n! \approx n \ln n - n + \tfrac12 \ln(2\pi n)$. Here $52 \ln 52 \approx 205.5$; subtract 52 and add $2.9$ to get $156.4$. Dividing by $\ln 10 \approx 2.303$ gives $67.9$, so $52! \approx 8 \times 10^{67}$. The rough form $n \ln n - n$ alone gives 66.6, which shows why the last term matters at this precision.

The Earth holds $6 \times 10^{27}$ grams, or $2 \times 10^{26}$ moles of atoms, or $1.2 \times 10^{50}$ atoms. The ratio is about $6.7 \times 10^{17}$. The machine has checked $10^9 \cdot 4.4 \times 10^{17} = 4.4 \times 10^{26}$ orderings, leaving a factor of $1.8 \times 10^{41}$.

For a match, use the birthday\index{birthday problem} estimate. With $k$ shuffles there are about $k^2/2$ pairs, each matching with chance $1/52!$, so an even chance needs $k \approx \sqrt{2 \ln 2 \cdot 52!} \approx 1.06 \times 10^{34}$.

\iqwhy{This is why randomised identifiers with enough bits are safe in practice: the square root in the birthday count is what sets the size, so 128 bits are comfortable where 64 are not.}
\end{solution}

\begin{solution}{iq:iv1:estimation:21}
Five or fewer has a chance of about $0.16\%$. Widen intervals by about $2.44$ times in logarithms, so ``within a factor of 3'' should become ``within a factor of about 15''.

The score is binomial with $n = 10$ and $p = 0.9$, and $P(X \le 5) \approx 0.00163$. So you are not unlucky: your intervals are too narrow. Even with so few, an exact 95\% interval for your real coverage runs from about 19\% to 81\%, and it does not reach 90\%.

For the scale, say your errors in logarithms are normal with true standard deviation $\sigma$. A real 90\% interval is $\pm 1.645\sigma$. You covered about 50\%, which is $\pm 0.674\sigma$. So what you called $1.645$ of your assumed $\sigma$ is really $0.674$ of the true one, and the ratio is $1.645/0.674 \approx 2.44$. A factor of 3 is $\ln 3$ in logs, and $e^{2.44 \ln 3} \approx 14.6$.

\iqwrong{Widening by the ratio of coverages, $90/50$. Coverage is not linear in width, because the tails of a normal are thin, so a good deal of extra width is needed to catch the last 40\%.}
\iqwhy{Overconfidence is the usual result on this kind of quiz, and interviewers who ask for an interval are watching for it. On a desk, this calibration is what sets position limits: being wrong outside your interval half of the time is a risk problem.}
\end{solution}

\begin{solution}{iq:iv1:estimation:22}
About $4.1\%$ if normal, and about 22\% under the Student $t$ with 3 degrees of freedom. A fall of about 20\% rejects the normal model, and it is not surprising under the fat-tailed one.

There are $25\,200$ days, so look for the move exceeded on one of them, which is a two-sided tail of $1/25\,200$. For a normal, $P(|Z| > z) = 1/25\,200$ gives $z \approx 4.11$, or $4.1\%$.

A $t$ with 3 degrees of freedom has variance 3, so it must be scaled by $1/\sqrt{3}$ to have a daily volatility of 1\%. Its tail falls only as a power\index{fat tails}, and the same tail point sits at about 22\%. The median of the 100-year maximum is a little higher, about $4.2\%$ and about 25\%, since it sits at a tail of $\ln 2/25\,200$.

A fall of 20 standard deviations under a normal has a chance of about $3 \times 10^{-89}$ a day, which is a wait of about $10^{86}$ years. Yet the Dow Jones Industrial Average fell $22.6\%$ in one session in October 1987, so that model must be the one to go.

\iqwhy{Estimates from a normal model look precise and can be wrong by a factor of five in the tail, where the capital is needed. Ask for the tail assumption whenever someone quotes a worst case.}
\end{solution}

\begin{solution}{iq:iv1:estimation:23}
The book can be \$$6.25$ billion. It makes about \$315 million a year, or about 5\% on the book.

The daily limit on trading is $1000 \cdot 50 \times 10^6 \cdot 5\% = \$2.5$ billion. Now count the trading carefully. To keep a book of size $G$ with a holding period of 5 days, each day you close $G/5$ and open another $G/5$, so you trade $2G/5$ a day. Setting $2G/5 = 2.5$ billion gives $G = 6.25$ billion.

The turnover is $2.5$ billion a day, or \$630 billion a year. At 5 basis points on each dollar traded, that is \$315 million, or $5.04\%$ on \$$6.25$ billion.

\iqwrong{Multiplying the daily limit by the holding period to get \$$12.5$ billion. That counts only the opening side. A position has to be closed too, and each of those trades uses the same volume budget.}
\iqfollow{Without changing the signal, what is the simplest way to raise capacity?}{Hold for longer, if the signal lasts. With a 10-day holding period the same trading supports twice the book, although the return per dollar of the book falls unless the signal grows in step.}
\iqwhy{Capacity is usually decided by this sum, not by the model. Investors ask for it early, because 5\% participation at five days is a statement about how large the fund can ever be.}
\end{solution}

\begin{solution}{iq:iv1:estimation:24}
The average wait is $1.6$ microseconds. The burst ends with a queue of 250 messages, and it empties about 333 microseconds later.

The handler's capacity is $1/0.8 = 1.25$ messages a microsecond, so its load is $\rho = 0.8$. With random arrivals and a fixed service time, the queue is an M/D/1 queue\index{M/D/1 queue}, and the Pollaczek-Khinchine formula gives
\[
W = \frac{\rho\, s}{2(1 - \rho)} = \frac{0.8 \cdot 0.8}{2 \cdot 0.2} = 1.6.
\]
Total time in the system is $2.4$ microseconds. If the service times were exponential the wait would double to $3.2$. The term $1/(1 - \rho)$ is the warning: at 90\% load the wait is more than double.

During the burst, the excess is $1.5 - 1.25 = 0.25$ messages a microsecond for 1\,000 microseconds, or 250 messages, which is a delay of 200 microseconds. Afterwards the queue shrinks at $1.25 - 0.5 = 0.75$ a microsecond, so it empties in $250/0.75 \approx 333$ microseconds. In other words, a burst of 1 ms hurts for about $1.3$ ms.

\iqwhy{Average load hides this risk. A handler at 40\% average load can still fall 200 microseconds behind in a burst, just as the market moves. Size for the burst, not for the average.}
\end{solution}

\begin{solution}{iq:iv1:estimation:25}
No: costs alone come to about 42 basis points per dollar traded, against about 6 of net profit claimed. With the 10-basis-point signal, the best book is about \$490 million, earning about \$83 million a year.

The fund trades $2 \cdot \$20$ billion a day, which is 8\% of the market, so $q = 0.08$, and the cost per dollar traded is $\sigma\sqrt{q} = 150 \cdot 0.283 \approx 42.4$ basis points. The claim is 30\% of \$20 billion, or \$6 billion a year. Spread over 252 days and \$40 billion a day, that is about $5.95$ basis points net. So the signal would have to earn about $48$ basis points on every dollar traded, every day, on a volume of about one-twelfth of the market. With a real signal of 10 basis points, this size loses about \$130 million a day.

Now optimise. Profit per day, with $q$ as the fund's share and $V$ the market's value, is
\[
\pi(q) = qV\bigl(a - \sigma\sqrt{q}\bigr),
\]
and setting the derivative to zero gives $a = \tfrac32\sigma\sqrt{q}$, which is $q^* = (2a/3\sigma)^2 \approx 0.00198$. So you trade about \$988 million a day, on a book of about \$494 million. A cost of two-thirds of the signal is the signature of the square-root law\index{square-root impact}. The net is $a/3 \approx 3.33$ basis points, or about \$329\,000 a day and \$83 million a year, about 17\% on the book.

\iqwrong{Stopping when costs equal the signal. Profit is then zero: the optimum is well before that, where the marginal dollar breaks even but the average dollar still earns.}
\iqwhy{This is the check to apply to any large claim, and this result is why good signals support small books: return falls as the book grows, while dollar profit rises much more slowly.}
\end{solution}
